\section{A retained prefix and its linked partition coefficient}
\label{sec:boundary}

\begin{proposition}[An actual $\omega^\omega$ block]\label{prop:omega-block}
Let $\Omega=\omega^\omega$. At a real suffix state $\pi+\Omega+n$, with $\pi$ a retained CNF
prefix (zero or of least exponent at least $\omega$), and total norm $u$,
assume $c=\norm(\pi)$, $c+3\le u$, and
$B=f(u)\ge\max(16,m)$. Put $C=f(u+1)$ and $d=B-c-1$.
If $H_c(u)$ counts the complete deletion of this $\Omega$ copy,
then
\begin{equation}\label{eq:omega-block}
 H_c(u)\equiv1+p_{B-c-1}(C-c)\pmod m.
\end{equation}
Its output norm is still exactly $f^{H_c(u)}(u)$.
\end{proposition}
\begin{proof}
The first predecessor is $\pi+\omega^d$, of total norm $B$.
Indeed the greatest exponent of norm at most $B-c-1$ is the
finite exponent $d$ in this particular bound below $\omega$;
this is the predecessor of $\omega^\omega$ under the suffix budget.
All subsequent norms $x$ are at least $B$. Their budgets
$G(x)=F_5(x+100)$ have residue $r_{Q_d}$ by the second part of
Theorem~\ref{thm:envelope}, because
\[
 Q_d=d!m^{d-1}\le B^{2B}\le2^{2^B}.
\]
The last inequality holds for $B\ge16$; it follows from
$2B\log_2 B\le2^B$, first at $16$ and then by the increasing
ratio of the right side to the left side.
The adjacent value $C=f(u+1)$ has the same residue modulo $Q_d$
by the third part of Theorem~\ref{thm:envelope}.
Apply Theorem~\ref{thm:transfer} to the remaining $\omega^d$ block
with $S=C$, then include its preceding $\Omega$ replacement.
Norm saturation proves the final assertion independently of the
modular calculation.
\end{proof}

\subsection{The first moving-prefix boundary}

Write $\Theta=\omega^{\omega+1}$, whose norm is four.
For each $v\ge4$, begin an actual descent at $\Theta+(v-4)$,
of total norm $v$. The same calculation applies whenever a larger
descent genuinely reaches this limit part at that norm; it does not
posit independent entry values inside a larger source.
Put $u=G(v)$. The residue recursion gives $u\equiv1\pmod3$.

\begin{lemma}[The retained prefix at first entry]\label{lem:entry}
The predecessor of this $\Theta$ state is
\[
 \Omega q+1,\qquad q=(u-1)/3.
\]
The first $\Omega$ copy deleted in this polynomial has retained
prefix norm $e=u-4$ and total input norm $u$.
\end{lemma}
\begin{proof}
A copy of $\Omega$ has norm three. The maximal affordable exponent
below $\omega+1$ is $\omega$, so the greedy predecessor takes as
many copies of $\Omega$ as possible. Since $u\equiv1\pmod3$,
exactly one norm unit remains, whose largest affordable ordinal
is the finite ordinal one. Before deleting the last displayed $\Omega$ in the
ordinary sum $\Omega(q-1)+(\Omega+1)$, the retained prefix is
$\Omega(q-1)$ and has norm $3(q-1)=u-4$.
\end{proof}

With $B=f(u)$, $C=f(u+1)$, $e=u-4$ and $d=B-u+3$, the
first block parity is consequently
\begin{equation}\label{eq:delta}
 \delta_4(v):=H_{u-4}(u)\bmod2
       =1+p_d(C-e)\pmod2.
\end{equation}
The subscript records the norm four of its suffix $\Omega+1$.
The initial $\Theta$ replacement lies outside $H_{u-4}(u)$;
the displayed $1$ instead counts the $\Omega$ replacement.

The value $C=G(u+1)$ is a modular surrogate. It is not the next
operational norm $G(B)$. The transfer theorem permits it because
its residue modulo the required partition period agrees with every
later budget. All four parameters $B,C,e,d$ are determined by the
single entry norm $v$.

This quantity is inside a larger correlated computation. Before its
remaining limit part has become exactly $\Omega$, the $\Theta$ segment
deletes $q-1$ copies with prefix masses $3(q-1),3(q-2),\ldots,3$. Their inputs
are successive points on the \emph{same} $f$ orbit. Observing the
sum of all these counts does not automatically define the graph of
the first summand. Thus \eqref{eq:delta} is a concrete benchmark for
a proposed arithmetic method, not a proved necessary subproblem of
every possible comparison procedure. The family
$\Theta+(v-4)$ proves genuine ordinal-descent realization for every
$v\ge4$; no stronger typed-source grammar realization is used here.

